\documentclass[11pt]{article}
\usepackage{mathblatt}
% Probe 08.10.2026, „Die lange Seite“ in drei Durchgängen: Übersicht übernommen aus paket-pythagoras-2, gilt für alle drei Durchgänge;
% Übersicht nach bau/bauregeln.md § 3 (Stand 08.10.2026).
% Bereiche als Überschrift, je Blatt eine Zeile, Name in Schülersprache; Fachwort nur, wenn es neu ist;
% kein grauer Zusatz; Rand-Blätter grau; das gelieferte Blatt ▸ fett; davor „Vorher“, danach „Weiter“.
\newcommand{\kennung}{H6Q \textperiodcentered{} H7Q \textperiodcentered{} H8Q}
\newcommand{\bereich}[1]{\par\addvspace{14pt}\noindent{\large\bfseries #1}\par\nobreak\vspace{4pt}}
\newcommand{\bl}[1]{\par\noindent #1\par\vspace{3pt}}
\newcommand{\blrand}[1]{\par\noindent{\color{mbgrau}#1}\par\vspace{3pt}}
\newcommand{\blhier}[1]{\par\noindent\llap{$\blacktriangleright$\hspace{0.5em}}\textbf{#1}\par\vspace{3pt}}
\begin{document}
\pfheftstil
\fancyfoot[C]{{\footnotesize\color{mbgrau}\kennung\hfill\thepage}}
\pfheftkopf{Satz des Pythagoras}{P10}

\bereich{Vorher}
\bl{Quadrat und Wurzel}
\bl{Rechter Winkel im Dreieck}

\bereich{Satz des Pythagoras}
\blhier{Die lange Seite \textperiodcentered{} Hypotenuse}
\blrand{Warum der Satz stimmt}
\bl{Eine kurze Seite \textperiodcentered{} Kathete}
\blrand{Ist der Winkel recht? \textperiodcentered{} Umkehrung}

\bereich{In Figuren und Körpern}
\bl{Das Dreieck erst finden}
\bl{Schräge Strecken im Körper}

\bereich{Weiter}
\bl{Seiten und Winkel mit sin, cos, tan}
\end{document}
